%% ===================== NEW SECTION: computational content =====================
\section{Computational content: a measured negative}
\label{sec:content}

The correlations above are large by the usual descriptive measures: $\delta = -0.01414$ with a
Pearson statistic of $10{,}167$ on one degree of freedom, and individual cross-base values of
$\phi/\sigma_0$ up to $476$. Those statistics are descriptive, not calibrated: they describe
$5 \times 10^7$ enumerated primes rather than a large effect. The question of this section is
whether the correlation is \emph{computationally} worth anything --- could it predict generator
availability, or yield a faster algorithm? We answer by measurement, and the answer is no. We
record it because the negative is informative and because the protocol transfers to any claimed
predictive advantage in this area.

\subsection*{The right baseline, and two different quantities}
A correlation can look large against a naive reference and carry no usable information once the
elementary structure is taken into account. For consecutive primes that structure is the
\emph{joint residue pair}: a practitioner asking whether $a$ is a primitive root mod $p$ already
knows $p \bmod 4\,\sqf(a)$, and the Legendre symbol $\Lsym{a}{p}$ is a function of it. The correct
baseline for a predictive claim is therefore not the marginal density of $\Art$ but the
conditional distribution given the joint residues.

Two quantities must be distinguished, and earlier drafts of this work conflated them.

\begin{itemize}
	\item An \emph{empirical conditional-entropy diagnostic}: estimate each cell's conditional
	probability from the data and score those same data. This measures how much of the
	variation the conditioning explains, and it is a legitimate descriptive summary --- but a
	finer nested model can never do worse, so it is \emph{not} evidence of predictive skill.
	\item A \emph{held-out} measurement: fit the cell probabilities on one part of the data and
	score a disjoint part. Only this can show predictive value.
\end{itemize}

Table~\ref{tab:content} gives the first quantity on all $50{,}847{,}530$ consecutive pairs below
$10^9$, and Table~\ref{tab:holdout} the second.

\begin{table}[t]
	\centering
	\begin{tabular}{lccc}
		\toprule
		conditioning & entropy (nats) & gain vs marginal & incremental \\
		\midrule
		marginal & $0.661047$ & --- & --- \\
		joint residues mod $120$ & $0.241529$ & $0.419517$ \ ($0.605$ bits) & --- \\
		\quad $+$ predecessor Artin status & $0.241520$ & --- & $\mathbf{0.000009}$ \\
		\midrule
		joint residues mod $840$ & $0.237235$ & $0.423812$ & --- \\
		\quad $+$ predecessor Artin status & $0.237217$ & --- & $\mathbf{0.000018}$ \\
		\bottomrule
	\end{tabular}
	\caption{\emph{Empirical conditional-entropy diagnostic} over $50{,}847{,}530$ pairs below
	$10^9$ (resubstitution: each cell's probability is estimated from the same pairs it scores).
	The elementary joint-residue baseline carries $0.42$ nats; the correlation adds of order
	$10^{-5}$ nats \emph{in this descriptive sense}. This is not a held-out measurement.}
	\label{tab:content}
\end{table}

\begin{table}[t]
	\centering
	\begin{tabular}{lccc}
		\toprule
		conditioning & held-out entropy (nats) & held-out gain \\
		\midrule
		joint residues mod $120$ & $0.241410$ & --- \\
		\quad $+$ predecessor Artin status & $0.241413$ & $-0.0000032$ \\
		\midrule
		joint residues mod $840$ & $0.237280$ & --- \\
		\quad $+$ predecessor Artin status & $0.237341$ & $-0.0000611$ \\
		\bottomrule
	\end{tabular}
	\caption{\emph{Held-out} measurement. Cell probabilities are fitted on pairs whose successor
	prime has an \emph{even} 64-bit hash (bit $0$ of the hash) and scored on the complementary
	odd-hash half ($25{,}421{,}106$ pairs);
	Jeffreys ($0.5$) smoothing is applied to empty and sparse cells. A hash split is used rather
	than a range split because the Artin density drifts with $x$, which would penalise the finer
	model for reasons unrelated to prediction. Both gains are \emph{negative}: adding the
	predecessor's status makes held-out prediction slightly worse.}
	\label{tab:holdout}
\end{table}

No estimator we tested converts the correlation into an improvement, but the in-sample
increment is not pure overfitting at modulus $120$, and we report both facts. The in-sample
increment is $0.000009$ nats (mod $120$) and $0.000018$ nats (mod $840$). The optimism expected
from the extra free parameters is the reference scale: splitting each cell on the predecessor's
status adds one probability per cell whose successor is not constant, and there are $k=256$
(mod $120$) and $k=1{,}702$ (mod $840$) such cells, so the classical $k/(2N)$ optimism is
$2.5\times10^{-6}$ and $1.7\times10^{-5}$ nats. At modulus $840$ the increment is therefore
indistinguishable from that scale; at modulus $120$ it exceeds it by a factor of about $3.6$.

A permutation test makes this concrete. Holding both margins of every cell fixed and redrawing
each $2\times2$ table from the hypergeometric distribution, the in-sample increment at modulus
$120$ lies far outside the null ($z=+26.9$ over $200$ replicates) while at modulus $840$ it does
not ($z=+2.4$). Applied to the held-out direction the same null predicts losses of
$-1.0\times10^{-5}$ and $-9.9\times10^{-6}$ at modulus $120$, against observed
$-3.2\times10^{-6}$ and $-4.9\times10^{-6}$ --- that is, the observed losses are much
\emph{smaller} than pure overfitting would produce ($z=+6.1$ and $+4.9$), which is positive
evidence of predecessor information. At modulus $840$ the held-out $z$ scores are $+1.6$ and
$+0.6$: no signal beyond the null.

The picture is therefore: a real but very small predecessor signal beyond the mod-$120$
residues, of order $10^{-6}$ nats by the split-sample estimate $(g_{\text{in}}+g_{\text{out}})/2$
($1.7\times10^{-6}$ and $8.5\times10^{-7}$ in the two directions); no signal detectable beyond
mod $840$; and in either case an edge several orders of magnitude below any threshold at which a
predictor could pay for itself. A pooled one-parameter variant (a single logit offset for the
predecessor's status on top of the residue cells) scores held-out gains of $-5\times10^{-8}$
(mod $120$) and $-8\times10^{-7}$ (mod $840$) --- no gain, within noise. The tension between
these two results is itself informative: a pooled offset buys nothing while cell-level
information is significant at modulus $120$, so the signal is heterogeneous across cells (a
sign-varying $\delta_c$) rather than a common shift. We state the negative as a model- and
workload-limited finding, not as an information-theoretic bound: nothing here bounds the oracle
conditional mutual information, other estimators or other ranges. One caution for the reader who
repeats this: conditioning on the gap class $g \bmod 40$ \emph{alone} produces a spurious
apparent gain of about $0.02$ nats, because $g \bmod 40$ does not determine the individual
residues $p_n \bmod 40$ and $p_{n+1} \bmod 40$; the apparent signal is residue information
re-appearing. The joint residues are the correct baseline.

\subsection*{What the between-cell share measures}
The global anticorrelation admits an exact two-term decomposition. Writing $X,Y$ for the
predecessor's and successor's Artin indicators and $C$ for the joint-residue cell, the law of
total covariance gives
\[
\delta \;=\; \underbrace{\frac{1}{\operatorname{Var}(X)}\sum_c \pi_c\, r_c(1-r_c)\,\delta_c}_{\text{within}}
      \;+\; \underbrace{\sum_c \bigl(w_c^{1}-w_c^{0}\bigr) p_c}_{\text{between}},
\]
with $\pi_c = P(C{=}c)$, $r_c = P(X{=}1\mid C{=}c)$, $w_c^{x}=P(C{=}c\mid X{=}x)$,
$p_c=P(Y{=}1\mid C{=}c)$ and $\delta_c$ the within-cell association. This is an identity, not a
model, and the driver accumulates both sums \emph{independently} and asserts it; at $10^9$ the
identity is reproduced to $2\times10^{-17}$ at both moduli. It gives $98.69\%$ (mod $120$) and
$99.55\%$ (mod $840$) of $\delta$ between joint-residue cells, with within-cell terms
$-0.000186$ and $-0.000064$.

The share depends on the weighting convention, and we state ours explicitly. Under the Kitagawa
ordering with $P(C{=}c\mid X{=}1)$ weights the between share is $96.1\%$ (mod $120$) and
$98.4\%$ (mod $840$); with $P(C{=}c\mid X{=}0)$ weights it is $99.2\%$ and $99.7\%$; with
symmetric weights $97.7\%$ and $99.1\%$. Every convention agrees that more than $96\%$ of the
anticorrelation at mod $120$, and more than $98\%$ at mod $840$, is attributable to the joint
residues rather than to association within them. Cells in which the predecessor's status never
varies carry no identifiable within-cell association; their mass (about half of all pairs at
mod $120$, where the residue class largely fixes Artin status) lands in the between term, which
is exactly where a deterministic exclusion belongs. The residual within-cell association,
$0.000186$ and $0.000064$, is the part a user holding the joint residues fixed would still have
to model.

\subsection*{Both exclusion laws are computationally redundant}
An exclusion law can be mathematically correct and computationally useless, and both of ours are
instances of that. The gap law (Theorem~\ref{thm:exclusion}) is the statement that for
$g \equiv 20 \pmod{40}$ the product $\chi_{10}(p)\chi_{10}(p+g)$ is the constant $\mathbf{-1}$
(equivalently $\chi_{10}\bigl(p(p+g)\bigr) = -1$), so the two characters are opposite and at most
one of the two primes can be Artin. The constant sign follows from quadratic reciprocity applied
to the shift $g$, together with multiplicativity; multiplicativity alone does not fix which sign
the shifted product takes. But a practitioner need not track predecessor status or gap state to
learn this: reading $p \bmod 40$ gives $\chi_{10}(p)$ directly, and $p+g \bmod 40$ gives
$\chi_{10}(p+g)$. The law is a statement about characters that a character computation already
answers.

The same is true of the same-prime law (Theorem~\ref{thm:pair}) and its triple corollary
(Theorem~\ref{thm:triple}): they follow from $\chi_c = \chi_a\chi_b$, so if $a$ and $b$ are Artin
then both characters are $-1$ and $\chi_c = +1$, whence $c$ is not Artin whenever the three bases are units modulo $p$
(i.e.\ $p\nmid abc$). To be precise about the scope of ``redundant'': \emph{once the relevant individual
quadratic characters are already available}, these statements supply no further rejection beyond
what those characters give. Both deductions combine characters that are already in hand; for
the gap law, reciprocity is what makes the $\chi_5$ character available in the first place.
Neither saves a
modular exponentiation that a character-aware implementation was not already saving. We verified
the gap law exhaustively at $10^9$ --- $2{,}214{,}511$ pairs with $g \equiv 20 \pmod{40}$, of
which $984{,}678$ have an Artin predecessor, and \emph{zero} are doubly-Artin --- but exhaustive
verification of a redundancy is still a redundancy.

\subsection*{The one usable filter is elementary}
The only filter that does real work is the necessary condition of \S\ref{sec:prelim}: a primitive
root must be a quadratic non-residue, so $\Lsym{a}{p}=-1$. It is answerable in $O(1)$ arithmetic
by reciprocity rather than by exponentiation, and it is exactly the information that the
conductor-$40$ quadratic channel of \S\ref{sec:decomposition} carries. That is the honest reading
of the decomposition: the paper \emph{explains} why a standard filter works, and identifies which
channel carries the effect, but it does not supply a new one.

We measured it in a scanning search for Artin-base-$10$ primes at $p \approx 10^{15}$, where
factorising $p-1$ dominates. The three strategies below are single-threaded loops on a 32-core
machine; the recorded times are kernel times for the search loop itself, with the shared
small-prime setup (identical for all three, and amortised in any real deployment) excluded.

\begin{table}[t]
	\centering
	\begin{tabular}{lccc}
		\toprule
		strategy & time (s) & factorizations of $p-1$ & ratio to naive \\
		\midrule
		naive (factor every prime) & $0.1642$ & $722$ & $1.00\times$ \\
		non-residue filter first & $0.0841$ & $376$ & $1.95\times$ \\
		\quad $+$ small-$\ell$ early exit & $0.0732$ & $300$ & $2.24\times$ \\
		\bottomrule
	\end{tabular}
	\caption{Finding $300$ Artin-base-$10$ primes above $10^{15}$; $722$ primes visited. The
	small-$\ell$ early exit tests the channels $\ell \in \{3,5,7,11,13,17,19\}$ of $p-1$ before
	the full factorisation, rejecting on first failure ($402$ cheap exponentiations). The last
	step is about $15\%$ faster, equivalently about $13\%$ less elapsed time, than the
	non-residue filter alone; the exact counters (primes, factorizations, exponentiations)
	reproduce independently, while timings are single-run and hardware-dependent.}
	\label{tab:filter}
\end{table}

Ordering the small prime channels of $p-1$ before the full factorisation --- testing
$a^{(p-1)/\ell} \not\equiv 1$ for small $\ell \mid p-1$ and rejecting on the first failure ---
saves $20\%$ of the factorizations. This is a \emph{test-ordering} principle rather than a new
algorithm, and it is workload-dependent in a way worth stating: if the task is to find \emph{any}
primitive root for a \emph{given} $p$, then $p-1$ is factored once and cached, and the early exit
saves the cheap exponentiations but not the factorisation at all. The gain exists only for
scanning searches over fresh primes.

\subsection*{What would change the verdict}
This is a negative assessment of the evidence, not an impossibility theorem. The conditioning
above does not prove exact conditional independence, and the verdict would change if: (i)~a
repeated workload existed in which many compatible primes or candidate bases are tested and the
required labels are already available; (ii)~a new \emph{deterministic} exclusion were found that
rules out a case the character deductions do not already rule out; (iii)~an actionable decision
followed from the correlation --- skipping a soundly excluded candidate, or reordering costly
tests --- with net cost falling after factorisation, feature acquisition and bookkeeping; and
(iv)~correctness were preserved, since a predictor may order tests but cannot replace
certification. None of these holds in the workloads tested here. We have not evaluated a
cross-base predictive workload at all, so the negative is demonstrated for the consecutive-prime
axis and for the scanning-search workload; the same-prime axis is covered only by the redundancy
argument above. The results are best understood as describing arithmetic structure, and as
identifying exactly which elementary structure the observed correlations are made of.
